\clearpage
\begingroup
\fancyhead[L]{\small THEORY OF THE MOTION OF THE MOON.}
\renewcommand{\headrulewidth}{0pt}
\clearpage
\SourceLeaf{617}
\begin{center}
\vspace*{.15\textheight}
{\Large\bfseries THEORY OF THE MOTION OF THE MOON.\par}
\end{center}
\vspace{.28\textheight}
\par\noindent\hfill 77\textsuperscript{*}\par

\clearpage
\SourceLeaf{619}
\begin{center}
{\large POSTHUMOUS PAPERS.\par}
\bigskip
{\Large\bfseries THEORY OF THE MOTION OF THE MOON.\par}
\bigskip
FIRST SECTION.
\medskip

\textit{Fundamental equations for the motion of the Moon.}
\bigskip

{\large Preliminary remarks.\par}
\end{center}

\par\noindent\hspace*{\parindent}Of the forces that influence the motion of the Moon, we consider here
only the gravitational attractions of the Earth, the Sun, and the Moon, which
we regard as concentrated at the centers of gravity of these celestial bodies. In
fact, the effects of the other planets are so extraordinarily small that
they entirely escape even the most delicate observations. The spheroidal
figure of the Earth, which modifies the latter assumption, likewise
gives rise to only a couple of equations amounting to a few seconds of arc. Moreover,
these circumstances can easily be taken into account as well, once the
motion of the Moon has been examined with these effects set aside.

\par\noindent\hspace*{\parindent}We shall also assume the motion of the Earth around the Sun to be elliptical
and set aside its perturbations by the Moon and the planets.
We therefore also regard the Earth's orbit around the Sun as lying
in a fixed plane (i.e. one that remains constantly parallel to itself),
and neglect both the gradual displacement
\SourceLeaf{620}of the plane itself, caused by the action of the planets,
and the periodic deviations of the Earth from it, which result from the
combined perturbations of the planets and the Moon. For all these
circumstances produce changes in the motion of the Moon around the Earth that
are imperceptible in themselves and may quite properly be left to a separate
investigation.

\par\noindent\hspace*{\parindent}To find the fundamental equations, we first refer the
positions of the Moon, Earth, and Sun in space to three fixed
planes that intersect at right angles at an entirely arbitrary point, or rather
at a point left unspecified, in space. We place one
of these planes parallel to the ecliptic and denote the distances
of the Moon, Earth, and Sun from it at any time \(t\) by
\(z,z',z''\), respectively. These distances are taken as positive on the side on which
the north pole of the ecliptic lies and as negative on the opposite side. We place the
second plane perpendicular to the first, so that the
line of intersection passes through the point on the celestial sphere where the vernal equinox
lies at a specified time \(T\). We denote the distances of the Moon, Earth, and
Sun from this plane, taken as positive on the side on which the
summer solstice point lies, by \(y,y',y''\). Finally, the
third plane is perpendicular to the two previous planes, and its line of intersection
with the first passes through the summer solstice point. We take the distances of the
Moon, Earth, and Sun from this plane as positive on the side of the
vernal equinox point and negative on the opposite side,
and denote them by \(x,x',x''\).

\par\noindent\hspace*{\parindent}Let the distances of the Earth from the Sun, of the Moon from the
Sun, and of the Moon from the Earth be \(R,R',R''\), respectively, and the masses of the Moon,
Earth, and Sun be \(m,m',m''\), respectively. The familiar principles of
dynamics then give the following six fundamental equations
\[
\begin{aligned}
0&=\frac{\mathrm{dd}x}{\mathrm{d}t^2}
  +\frac{m'(x-x')}{R''^3}+\frac{m''(x-x'')}{R'^3}\\[4pt]
0&=\frac{\mathrm{dd}y}{\mathrm{d}t^2}
  +\frac{m'(y-y')}{R''^3}+\frac{m''(x-x'')}{R'^3}\\[4pt]
0&=\frac{\mathrm{dd}z}{\mathrm{d}t^2}
  +\frac{m'(z-z')}{R''^3}+\frac{m''(z-z'')}{R'^3}
\end{aligned}
\]
\SourceLeaf{621}
\[
\begin{aligned}
0&=\frac{\mathrm{dd}x'}{\mathrm{d}t^2}
  +\frac{m(x'-x)}{R''^3}+\frac{m''(x'-x'')}{R^3}\\[4pt]
0&=\frac{\mathrm{dd}y'}{\mathrm{d}t^2}
  +\frac{m(y'-y)}{R''^3}+\frac{m''(y'-y'')}{R^3}\\[4pt]
0&=\frac{\mathrm{dd}z'}{\mathrm{d}t^2}
  +\frac{m(z'-z)}{R''^3}+\frac{m''(z'-z'')}{R^3}
  \mathpunct{\raisebox{.4ex}{.}}
\end{aligned}
\]
\clearpage
\endgroup
